% =============================================================================
% Project : THE PHILOSOPHY OF VIABILITY
% Book    : Book I — Ontology of Viability
% Book ID : VS.B01
% Asset   : TL01 — Leibnizian Predicate Completeness and Ontological Identity
% Asset ID: VS.B01.P01.C08.S07.SS08.P006.TL01
% Path    : 4_ASSETS/tables/
% =============================================================================

\begin{table}[htbp]
\centering
\small
\caption{Predicate completeness and ontological identity in the Leibnizian comparison}
\label{tab:leibniz-predicate-completeness}
\renewcommand{\arraystretch}{1.2}
\begin{tabular}{|p{0.20\textwidth}|p{0.35\textwidth}|p{0.35\textwidth}|}
\hline
\textbf{Question} & \textbf{Leibnizian commitment} & \textbf{Comparative question from Part I} \\
\hline
Individual concept &
The complete concept contains every predicate that is truly attributable to the individual. &
Does the complete determination of an individual exhaustively explain its ontological identity? \\
\hline
Identity of indiscernibles &
No two distinct individual substances differ only numerically while sharing all their relevant determinations. &
Does the absence of discernible difference establish the identity of the individuals, or only constrain their possible distinction? \\
\hline
Determination and identity &
The individual is characterized through the complete set of predicates belonging to its concept. &
Part I distinguishes Identity from Determination; the relevance of determinations does not by itself establish their identity with Identity. \\
\hline
Unity and individuality &
A monad is a simple substance and a genuine unity within Leibniz's metaphysics. &
Unity and simplicity do not, without further argument, provide an exhaustive account of Identity. \\
\hline
Expression &
Each monad expresses the universe according to its own determinate manner. &
Expression must not be identified with Identity or transferred without qualification into the distinction between Expression and Manifestation in Part I. \\
\hline
Persistence and change &
The monad's succession of internal states is explained within Leibniz's account of simple substance. &
This account does not automatically settle the broader questions of persistence, continuity, and transformation raised in Part I. \\
\hline
\end{tabular}
\end{table}

% =============================================================================
% END OF ASSET TL01
% =============================================================================
